\documentclass[11pt]{article}

\usepackage[margin=1in]{geometry}
\usepackage{amsmath,amssymb,amsthm}
\usepackage{enumitem}

\newtheorem{example}{Example}
\newtheorem{corollary}{Corollary}

\newcommand{\WSPT}{\mathrm{WSPT}}
\newcommand{\OPT}{\mathrm{OPT}}
\newcommand{\Cminmax}{C_{min}^{max}}

\begin{document}

Now, we introduce a transformation that replaces a set $M$ of positive
numbers with a number $p$ that is used $t$ times such that the following
equations hold:
\begin{align}
        \sum_{i\in M} i &= t\cdot p, \\
        \sum_{i\in M} i^2 &= t\cdot p^2 .
\end{align}
Eqs.~(4) and (5) lead to a unique solution with
\[
        t=\frac{\left(\sum_{i\in M}i\right)^2}{\sum_{i\in M}i^2}
        \qquad\text{and}\qquad
        p=\frac{\sum_{i\in M}i^2}{\sum_{i\in M}i}.
\]
In Eq.~(2), this transformation is either applied to the set $L(I)$ of long
jobs or to some makespans. However, $t$ need not be an integer but may be a
rational number with denominator $q$ resulting in an instance with a
fractional number of jobs. To obtain an instance with an integral number of
jobs, we simply replace each machine and each job of the original instance
$I$ with $q$ identical copies. In the \(\mathrm{WSPT}\) order, all replicas of a job occur
one after another. Therefore, the schedule of each machine is replicated for
schedule $S_{\WSPT}$. Replicating the schedule of each machine for schedule
$S_{\OPT}$ also yields an optimal schedule for the same reasons discussed
before. Hence, $C(I)(\alpha)$ and $C^*(I)(\alpha)$ are simply multiplied by $q$, and
$\lambda(I)$ remains unchanged.

The application of the transformation to the set $L(I)$ or to machine
makespans of a list schedule $\sigma(I)$ does not change
$\sum_{i=1}^m C_{\max}^i(\sigma(I))$ due to (hint to Corr. 3) Furthermore, due to the
convexity of the square function, we have
\begin{equation}
\begin{aligned}
        t
        &=
        \frac{\left(\sum_{j\in L(I)}p_j\right)^2}
        {\sum_{j\in L(I)}p_j^2}                                      \\
        &\leq
        \frac{\left(\sum_{j\in L(I)}p_j\right)^2}
        {|L(I)|\left(\frac{\sum_{j\in L(I)}p_j}{|L(I)|}\right)^2}       \\
        &=
        \frac{|L(I)|\left(\sum_{j\in L(I)}p_j\right)^2}
        {\left(\sum_{j\in L(I)}p_j\right)^2}
        =
        |L(I)|.
\end{aligned}
\end{equation}


\begin{corollary}
For every instance $I$ of the $P||\sum_j p_j C_j(\alpha)$ problem, there is an
instance $I'$ of the $P||\sum_j p_j C_j(\alpha)$ problem with
$\lambda(I')(\alpha)\geq\lambda(I)(\alpha) \, \forall \alpha \in [\frac12,1]$ such that instance $I'$ contains at most $m-1$
jobs with the same processing time while all its other jobs have an
arbitrarily small processing time.
\end{corollary}

\begin{proof}
Due to Corollary 3, assume schedule $S_{\WSPT}(I)$ for instance $I$ in which
only jobs with an arbitrarily small processing time start before
$\Cminmax(S_{\WSPT}(I))$.

We apply our transformation to the set $L(I)$ to generate an instance $I'$ in
which all $t$ long jobs have the same processing time. In the \(\mathrm{WSPT}\) order of
both instances, long jobs appear after all small jobs. Therefore, we have
\[
        \Cminmax(S_{\WSPT}(I))=\Cminmax(S_{\WSPT}(I')).
\]
We discuss the consequences of the transformation to the primary \(\mathrm{WSPT}\) schedule
and an optimal schedule separately. Due to Eqs.~(3)--(5), we have
\[
\begin{aligned}
        2C(I)(\alpha)
        &=
        (2\alpha-1)\sum_{j\in L(I)}p_j^2
        +
        \sum_{j\in L(I)}(p_j+\Cminmax(S_{\WSPT}(I)))^2                  \\
        &\quad
        +(m-|L(I)|)(\Cminmax(S_{\WSPT}(I)))^2                           \\
        &=
        (2\alpha-1)\sum_{j\in L(I)}p_j^2
        +
        \sum_{j\in L(I)}p_j^2
        +
        2\sum_{j\in L(I)}(p_j\Cminmax(S_{\WSPT}(I)))               \\
        &\quad
        +m(\Cminmax(S_{\WSPT}(I)))^2                                    \\
        &=
        (2\alpha-1)tp^2+tp^2
        +2tp\Cminmax(S_{\WSPT}(I'))
        +m(\Cminmax(S_{\WSPT}(I')))^2                                  \\
        &=
        (2\alpha -1)tp^2
        +t(p+\Cminmax(S_{\WSPT}(I')))^2
        +(m-t)(\Cminmax(S_{\WSPT}(I')))^2                              \\
        &=
        2C(I')(\alpha).
\end{aligned}
\]

Next we address the optimal schedules for instances $I$ and $I'$. To show
that $C^*(I')\leq C^*(I)$ holds, we consider two cases:

\begin{enumerate}[label=(\arabic*)]
\item
$C_{\max}^i(S_{\OPT}(I'))=\Cminmax(S_{\OPT}(I'))$ holds for all machines $i$.

Remember that any list schedule $\sigma(I')$ is optimal if
$C_{\max}^i(\sigma(I'))=\Cminmax(\sigma(I'))$ holds for all machines $i$ due to the
convexity of the square function. Furthermore, the transformation does not
change $\sum_{i=1}^m C_{\max}^i(\sigma(I))$ for any list schedule $\sigma(I)$.
Therefore, eqs. whaterber is needed here yield together with the convexity of the square
function:
\[
\begin{aligned}
        2C^*(I)(\alpha)
        &=
        (2\alpha-1)\sum_{j\in L(I)}p_j^2
        +
        \sum_{i=1}^m(C_{\max}^i(S_{\OPT}(I)))^2                 \\
        &=
        (2\alpha-1)t\cdot p^2
        +
        \sum_{i=1}^m(C_{\max}^i(S_{\OPT}(I)))^2                 \\
        &\geq
        (2\alpha-1)t\cdot p^2
        +
        m(\Cminmax(S_{\OPT}(I')))^2
        =
        2C^*(I')(\alpha).
\end{aligned}
\]

\item
$C_{\max}^i(S_{\OPT}(I'))>\Cminmax(S_{\OPT}(I'))$ holds for some machines $i$.
In the following sequence of equations, we show that $C^*(I')\leq C^*(I)$
also holds for this case. The individual equations are later explained in
detail.
\begin{align}
        2C^*(I)(\alpha)
        &=
        (2\alpha-1)\sum_{j\in L(I)}p_j^2
        +
        \sum_{j\in K(I)}p_j^2
        +
        (m-|K(I)|)(\Cminmax(S_{\OPT}(I)))^2                      \tag{7}\\
        &=
        (2\alpha-1)t p^2+t'(p')^2
        +(m-|K(I)|-k)(\Cminmax(S_{\OPT}(I)))^2                   \tag{8}\\
        &\geq
        (2\alpha-1)tp^2+t'(p')^2
        +(m-t')
        \left(
            \frac{m-|K(I)|-k}{m-t'}\Cminmax(S_{\OPT}(I))
        \right)^2                                                \tag{9}\\
        &\geq
        (2\alpha-1)tp^2+tp^2
        +(m-t)(\Cminmax(S_{\OPT}(I')))^2                         \\
        &=
        2\alpha tp^2+(m-t)(\Cminmax(S_{\OPT}(I')))^2
        =
        2C^*(I')(\alpha).                                                \tag{10}
\end{align}

First, we introduce the size
\[
        k=
        \frac{\sum_{j\in L(I)\setminus K(I)}p_j}
        {\Cminmax(S_{\OPT}(I))}
        \leq
        |L(I)|-|K(I)|.
\]
If necessary, we use replication to guarantee an integral $k$. Eq.~(7) is
the application of Eq.~(2) on schedule $S_{\OPT}(I)$ by using set $K(I)$.
We apply our transformation twice to Eq.~(7):
\begin{enumerate}[label=(\alph*)]
\item on set $L(I)$ to transform $\sum_{j\in L(I)}p_j^2$;
\item on the makespan of all machines that execute jobs in $K(I)$ and of
$k$ additional machines with
$C_{\max}^i(S_{\OPT}(I))=\Cminmax(S_{\OPT}(I))$.
\end{enumerate}
The second transformation produces the values $t'$ and $p'$. Remember that
we use replication to guarantee an integral $t'$. Due to the convexity of the
square function and $0\leq p_j\leq \Cminmax(S_{\OPT}(I))$ for all jobs
$j\in L(I)\setminus K(I)$, we have
\begin{align}
        t'\cdot p'
        &=
        \sum_{j\in K(I)}p_j+k\cdot\Cminmax(S_{\OPT}(I))
        =
        \sum_{j\in L(I)}p_j
        =
        t\cdot p,                                                \tag{11}\\
        t'\cdot(p')^2
        &=
        \sum_{j\in K(I)}p_j^2+k(\Cminmax(S_{\OPT}(I)))^2          \notag\\
        &=
        \sum_{j\in K(I)}p_j^2+k(\Cminmax(S_{\OPT}(I)))^2
        +(|L(I)|-|K(I)|-k)\cdot 0^2                              \notag\\
        &\geq
        \sum_{j\in K(I)}p_j^2
        +
        \sum_{j\in L(I)\setminus K(I)}p_j^2
        =
        \sum_{j\in L(I)}p_j^2
        =
        t\cdot p^2,                                              \tag{12}\\
        &\Rightarrow
        p'\geq p
        \quad\text{and}\quad
        t'\leq t.                                                \tag{13}
\end{align}

In Eq.~(8), the result of the second transformation to Eq.~(7) is shown.
Note that both transformations are formally applied and do not produce a
valid schedule. In particular, $|K(I)|+k-t'$ machines are not used in
Eq.~(8) as $t'\leq |K(I)|+k$ holds due to Eq.~(6).

In Eq.~(9), we evenly distribute the total load
\[
        (m-|K(I)|-k)\Cminmax(S_{\OPT}(I))
\]
of all small jobs on $m-t'$ machines. This distribution reduces the total
weighted completion time due to the convexity of the square function.

Eq.~(10) shows schedule $S_{\OPT}(I')$. In this schedule, only
$m-t\leq m-t'$ machines are available to execute the small jobs due to
Eq.~(13). Therefore, we have
\[
        \frac{m-|K(I)|-k}{m-t'}\Cminmax(S_{\OPT}(I))
        \leq
        C_{\max}^i(S_{\OPT}(I')).
\]
Due to this relation, Eqs.~(11)--(13), and the convexity of the square
function, the total weighted completion time is reduced again.
\end{enumerate}
\end{proof}

\end{document}
